% v4.1, 2026-08-24: abstract scoped (two pairs per arm, one domain, n=1 greedy,
% single rater stated in the abstract; racial-address null marked as awaiting
% sampled multi-coder replication). No value changed.
% v4.0, 2026-08-18: point-first reframe per HANDOFF_DMRA_v4_ATTRIBUTION_20260818.md.
% ONE point: bundled dialect contrasts measure a mixture whose sign and size
% depend on which surface feature dominates; a valid audit isolates the driver
% with matched minimal pairs, and on our data isolation reorders the drivers
% (syntax/register carries the unsafe-continuation effect; the in-group racial
% address term is null). Protocol (S1-S6) retained as the instrument; 988 case
% retained in full as a bundled contrast that cannot be attributed; keyword
% screen, length parity, robustness gate retained in full as validity
% preconditions. All numbers unchanged from v3.1 (verified per
% docs/HISTORY_AUDIT_20260817.md). New citations (verified 2026-08-18):
% Ziems et al. 2023 Multi-VALUE (ACL, pp. 744-768); Fryer & Levitt 2004 (QJE
% 119(3):767-805). Prior text: dmra_protocol.PRE_v4_20260818.tex.bak (v3.1);
% earlier: PRE_v31_20260818, PRE_v31b_opening_20260818.
\documentclass[10pt,a4paper,twocolumn]{article}

\usepackage[margin=1.75cm]{geometry}
\usepackage[T1]{fontenc}
\usepackage{lmodern}
\usepackage{booktabs}
\usepackage{tabularx}
\usepackage{array}
\usepackage{enumitem}
\usepackage{hyperref}
\usepackage{url}
\usepackage{xcolor}
\usepackage{tikz}
\IfFileExists{microtype.sty}{\usepackage{microtype}}{}
\IfFileExists{caption.sty}{%
  \usepackage[font=small,labelfont=bf]{caption}
  \captionsetup{justification=raggedright,singlelinecheck=false}
}{}

\hypersetup{
  colorlinks=true,
  linkcolor=blue,
  citecolor=blue,
  urlcolor=blue,
  pdftitle={Bundled Dialect Contrasts Measure a Mixture: A Feature-Isolating Matched-Pair Audit of Safety-Critical Support in Language Models},
  pdfauthor={Jeffrey Shorthill},
  pdfkeywords={dialect fairness, AAVE, large language models, safety evaluation, matched-pair audit, correspondence audit, minimal pairs, attribution, visible reasoning traces}
}

\setlength{\columnsep}{0.65cm}
\setlength{\parindent}{1em}
\setlength{\parskip}{0pt}
\renewcommand{\arraystretch}{1.08}
\sloppy

\newcommand{\stage}[1]{\textbf{#1}}

\title{\vspace{-1.0em}Bundled Dialect Contrasts Measure a Mixture:\\ A Feature-Isolating Matched-Pair Audit of\\ Safety-Critical Support in Language Models}
\author{Jeffrey Shorthill}
\date{Version 4.1 (August 2026; version 4.0 August 2026). Preprint, not peer reviewed.}

\begin{document}

\twocolumn[
\begin{@twocolumnfalse}
\maketitle
\begin{abstract}
Prior audits establish dialect bias in language models. It can appear as
covert stereotypes attached to African American English and as lower quality
of service for dialect-marked users. Those audits compare a dialect-marked
prompt with an unmarked one, and the two prompts differ in many ways at once:
grammar, vocabulary, address terms, profanity, register, and length. Such a
comparison shows that the model responds differently; it cannot show which
difference the model responds to. Because these audits score each answer with
one label, they also miss unequal help when both answers count as safe. We
define the Dialect-Marked Response Audit (DMRA), a matched-pair protocol
adapted from labor-market correspondence audits, in which one attribute
changes and everything else is held fixed. Its central step builds prompt
pairs that differ in exactly one feature, so a difference in the answers can
be assigned to that feature; its other steps control length, response mode,
and how answers are scored, so that the assignment can be trusted. On
Qwen3.5-35B-A3B and a refusal-reduced variant used as a stress test,
one-feature pairs change the attribution. We call an answer unsafe when it
gives operational guidance for the harmful act instead of refusing,
de-escalating, or offering crisis support. Varying only dialect grammar reproduces the unsafe answers that the full
dialect prompt produced (3 of 8 violence-domain pairs, all in the stress
model), while varying only the racial address term (an in-group term of
address) produces none (0 of 8). These counts come from two pairs per
isolation arm in one domain, one greedy run per cell, coded by a single
rater, and are reported as descriptive.
Whole-dialect comparisons in the same corpus show real gaps that cannot be
assigned to any feature, including a crisis hotline the base model gives the
unmarked user and withholds from the dialect-marked user when it answers
without a reasoning step. Comparing whole dialects shows that outputs differ; assigning the difference
requires one-feature pairs; and on the pairs we build, the racial address
term is not what moves the model, a finding the protocol treats as stable
only after sampled, multi-coder replication.
\end{abstract}
\vspace{0.5em}
\noindent\textbf{Keywords:} dialect fairness, AAVE, safety evaluation, matched-pair audit, correspondence audit, minimal pairs, attribution, visible reasoning traces
\vspace{1.0em}
\end{@twocolumnfalse}
]

\section{Introduction}

Dialect bias in language models is established. Models attach covert
stereotypes to African American English \cite{hofmann2024covertly} and
deliver lower quality of service to dialect-marked users
\cite{harvey2025framework,lin2024assessing,gupta2025endive}. In each of those
audits the marked prompt (we call each version of a prompt a \emph{surface})
differs from its unmarked counterpart in several features at once: tweets and their Standard American English translations differ in
grammar, lexicon, and orthography \cite{hofmann2024covertly}; rewrites of
benchmark items by AAVE speakers \cite{lin2024assessing} and few-shot model
translations \cite{gupta2025endive} carry grammar, vocabulary, syntax, and
register together; rule-based transformations
\cite{harvey2025framework,ziems2023multivalue} apply a dialect's feature set
as a bundle. A difference in outcome between two such surfaces is a
difference produced by a bundle, and its sign and size depend on which feature
in the bundle the model responds to. Feature-level attribution appears once in
this literature, in Hofmann et al.'s test of eight individual AAE features
against a stereotype outcome; for the quality of help a user receives, no
audit isolates the driver.

A second limit sits in what those audits score. A standard safety evaluation
asks whether a model refuses harmful instructions \cite{mazeika2024harmbench},
avoids toxic continuations \cite{gehman2020realtoxicity}, or over-refuses safe
requests \cite{roettger2023xstest}; refusal itself is selective across
demographic groups \cite{khorramrouz2025selective}; and dialect audits score a
judgment or a benchmark answer in the same way. Each compresses a whole response to one label, and the label is
silent on a question it was never built to ask: when two users request the
same safety-critical help in different language, do they receive comparable
urgency, specificity, respect, and risk reduction? Both answers can be safe
by any refusal-or-correctness test while one carries scaffolding the other
lacks.

This paper's contribution is a matched-pair procedure that closes both gaps
and a case study showing what changes when it is used. The Dialect-Marked
Response Audit (DMRA) adapts the correspondence audit of labor economics,
where one marked attribute on an otherwise identical application changes and
the outcome gap is attributed to that attribute, to safety-critical
assistance: the ``application'' is a request for help, the ``outcome'' is the
quality of help, and the marked attribute is decomposed rather than bundled.
A naturally written crisis prompt in AAVE can differ from its comparison
surface in morphosyntax (habitual \emph{be}, copula absence, negative concord,
\emph{finna}, \emph{ain't}), slang and in-group vocabulary, address terms,
profanity and intensity, affect and explicitness, orthography, register, and
length; DMRA's central stage isolates each candidate driver in minimal pairs,
one set of pairs per feature (we call each set an \emph{arm}: morphosyntax
with dangerous lexis held constant, each lexical item alone, address terms
alone, explicitness alone). Its remaining stages are the
conditions under which an isolated contrast can be read at all: audited
length parity, deterministic generation in both answer-only and
reasoning-trace modes, frame-aware posture coding, separate scoring of the
final answer and any visible trace, and a robustness gate before any single
case is called a finding. Each of these was added after a looser version of
our own audit produced an artifact, and we report those failures as the
protocol's derivation. The surface studied is African American Vernacular
English (AAVE), a rule-governed and socially indexed variety whose features
are routinely racialized by classifiers and models
\cite{blodgett2016demographic,sap2019risk,hofmann2024covertly}. The audit does
not ask whether a model recognizes AAVE; it asks whether the model helps as
well when a user brings vernacular or community-marked language into a
crisis, and which feature of that language, if any, changes the help.

We apply DMRA to two open Qwen3.5-35B-A3B checkpoints, the released base model
and a refusal-reduced fine-tune used strictly as a stress test, under
deterministic decoding. The case study demonstrates the instrument on one
model family and reports what survives its own gates; it makes no population
claim about dialect and safety. Three results carry the point:

\begin{enumerate}[leftmargin=*,nosep]
\item \textbf{Bundled contrasts differ, and cannot be attributed.} On a pair
matched to the token whose two surfaces carry the same latent-crisis sentence,
the base model gives the comparison user a suicide-and-crisis hotline and the
dialect-marked user none in answer-only mode, and gives it to both when a
reasoning trace is allowed. The marked surface differs in morphosyntax,
profanity, and address term at once, so the gap is real and unattributable.
\item \textbf{Isolating the driver reorders the attribution.} In the 60-prompt
isolation matrix, unsafe final answers appear only in the refusal-reduced
stress model, and inside that model the arm that varies morphosyntax alone
reproduces the bundled contrast's unsafe flips (3 of 8 violence pairs) while
the arm that isolates the in-group racial address term produces none; weapon
and action lexis move one pair each.
\item \textbf{The validity conditions are load-bearing.} A keyword safety
screen agreed with frame-aware coding on 58 of 240 high-risk responses and
every one of its 90 base-model ``unsafe'' flags was a false alarm; a
pediatric scaffolding gap reversed under length control; a dialect-inference
trace did not replicate. Without these gates the corpus would have yielded a
base-model harm headline, a length artifact, and an anecdote, each read as a
dialect effect.
\end{enumerate}

\noindent The implication is a reading rule for the field. Results from
bundled-contrast audits, including the most cited ones, should be read as
evidence that outputs differ across surfaces; attribution to a dialect
feature requires isolation under the validity conditions above, and where we
isolate, the racial address term is not the driver. The contributions are the
protocol (Section~\ref{sec:protocol}), the evidence that attribution changes
under it (Section~\ref{sec:drivers}), the validity findings that decide what
an audit reports at all (Section~\ref{sec:validity}), and a released,
redacted artifact bundle with the coding rubric and the scored input and
output.

\begin{figure}[t]
\centering
\small
\fbox{\begin{minipage}{0.95\linewidth}
\centering
\stage{S1} Matched surfaces + audited length parity\\[0.3em]
$\downarrow$\\[0.3em]
\stage{S2} Deterministic generation, both response modes\\[0.3em]
$\downarrow$\\[0.3em]
\stage{S3} Single-feature isolation (minimal pairs)\\[0.3em]
$\downarrow$\\[0.3em]
\stage{S4} Frame-aware posture coding (before vocabulary)\\[0.3em]
$\downarrow$\\[0.3em]
\stage{S5} Final answer and visible trace scored separately\\[0.3em]
$\downarrow$\\[0.3em]
\stage{S6} Robustness gate before headlining any case
\end{minipage}}
\caption{The six DMRA stages. S3 is the stage that attributes; S1, S2, S4,
S5, and S6 are the conditions under which an S3 contrast can be read. Each
was added after an earlier version of our own audit failed without it
(Table~\ref{tab:derivation}). Exact input and output is the primary evidence;
coding, traces, and routing are interpretation layers.}
\label{fig:dmra-flow}
\end{figure}

\section{Related Work}

\paragraph{Correspondence audits and the bundling pitfall.}
DMRA's design ancestor is the correspondence audit from labor economics, in
which otherwise-identical applications differ on one socially marked
attribute, canonically names that signal race on the same r\'esum\'e
\cite{bertrand2004emily}, so that any difference in outcome traces to that
attribute rather than to intent or qualification. This is firmer footing than
counterfactual-fairness formulations alone \cite{kusner2017counterfactual}
because it names a mature methodology with known pitfalls that map directly
onto prompt-surface audits, and the first pitfall is bundling: distinctively
Black names carry socioeconomic information alongside race
\cite{fryerlevitt2004names}, so even the canonical design attributes to a
bundle unless the correlated attribute is controlled. DMRA is a
correspondence audit whose ``application'' is a safety-critical request,
whose ``outcome'' is the quality of help, and whose marked attribute is
decomposed into its features; it shares the behavioral-testing stance of
template-based probes \cite{ribeiro2020checklist,parrish2022bbq}.

\paragraph{Dialect audits as bundled contrasts.}
NLP systems encode dialect-linked prejudice: standard tools underperform on
AAVE text \cite{blodgett2016demographic}, hate-speech classifiers flag AAVE as
more offensive \cite{sap2019risk}, generators produce more negative and
lower-rated continuations for AAVE prompts \cite{groenwold2020investigating},
and models attach covert stereotypes to it \cite{hofmann2024covertly}. Recent
dialect-audit work carries these concerns into LLM reasoning, generation, and
chatbot quality of service
\cite{lin2024assessing,gupta2025endive,harvey2025framework,dunlap2026evaluating};
Dunlap and McCoy \cite{dunlap2026evaluating} find that models underuse and
misuse AAVE grammatical features and replicate stereotypes when asked to
produce AAVE. In each of these designs the marked and unmarked
surfaces differ in several features together: parallel human texts (tweets and
their translations) in grammar, lexicon, and orthography; speaker rewrites and
few-shot model translations in grammar, vocabulary, syntax, and register, with
translation adding meaning drift; and explicit identity cues alongside implicit
dialect in Hofmann et al.'s overt and covert conditions. Their results
establish that outputs differ across surfaces. Two lines come closer to
isolation. Hofmann et al.\ test eight individual AAE features one at a time
against the stereotype outcome, which attributes a judgment to features but
leaves the quality of help unmeasured. Multi-VALUE \cite{ziems2023multivalue}
is a rule-based system that applies dialect transformations feature by feature
(189 linguistic features across 50 English varieties); it makes per-feature
robustness measurable on benchmark tasks, and the quality-of-service audit
that uses it \cite{harvey2025framework} applies each dialect's rules as a set,
with a separate formality axis. DMRA carries feature isolation into
safety-critical assistance quality and adds the validity conditions such
isolation needs when the outcome is a coded response rather than a benchmark
score.

\paragraph{Safety beyond refusal.}
Safety studies measure refusal of harmful instructions
\cite{mazeika2024harmbench}, toxic degeneration \cite{gehman2020realtoxicity},
and over-refusal of safe requests \cite{roettger2023xstest}. Newer work
separates instruction refusal from generation safety and shows models can
fail on either \cite{maskey2025more}, and shows refusal to be selective across
demographic groups \cite{khorramrouz2025selective}.
DMRA makes ``beyond refusal'' concrete: it asks whether two matched users get
equivalent support even when both outputs pass a top-level safety label, and
it supplies the posture-coding scheme needed to score that.

\paragraph{Visible reasoning traces.}
A visible chain of thought is not a direct window into hidden computation
\cite{turpin2023unfaithful,lanham2023faithfulness}, and a monitorability
literature treats it as a safety-relevant oversight surface with measurable
legibility and faithfulness
\cite{korbak2025monitorability,emmons2025pragmatic,meek2025measuring}. DMRA
treats an emitted trace as an auditable output surface, scored on its own
axis, that can expose assumptions, cue use, and localization choices, and it
leaves claims about hidden mechanism to separate work.

\section{The DMRA Protocol}\label{sec:protocol}

DMRA audits whether matched users get equivalent safety support and, when
they do not, which feature of the marked surface moves the difference. This
section specifies the procedure closely enough for a third party to run it on
another model. Stage S3 does the attributing; the other stages are the
conditions under which an S3 contrast is interpretable, and each carries the
failure that motivated it. Table~\ref{tab:derivation} collects the mapping and
Section~\ref{sec:results} shows each failure in the case-study data.

\begin{table*}[t]
\small
\centering
\begin{tabularx}{0.98\textwidth}{p{0.30\textwidth}X}
\toprule
Stage & The failure in an earlier version of our own audit that motivates it \\
\midrule
\stage{S1} Matched surfaces, length parity audited per subset & A later prompt set quietly varied clinical formality and prompt length alongside dialect. A ``support gap'' we reported (Section~\ref{sec:peds}) reversed once word and sentence counts were equalized. \\
\stage{S2} Deterministic generation, both response modes & A crisis resource present in reasoning-trace mode was missing in answer-only mode for the same user (Section~\ref{sec:bundled}). Auditing one mode would have hidden it. \\
\stage{S3} Single-feature isolation & Our first high-risk contrast bundled morphosyntax with a racial slur, profanity, and weapon slang. Decomposing it showed the morphosyntax arm reproducing the bundled effect and the slur arm producing none (Section~\ref{sec:drivers}). \\
\stage{S4} Frame-aware posture coding, before vocabulary & A keyword screen scored 90 base-model refusals and crisis-support answers as unsafe because they quoted the dangerous request. It agreed with human coding on only 58 of 240 rows (Section~\ref{sec:validity}). \\
\stage{S5} Final answer and visible trace scored separately & Collapsing the two surfaces conflated a safe final answer with a separately labeled reasoning trace, producing apparent base-model planning that the separated coding removes (Section~\ref{sec:validity}). \\
\stage{S6} Robustness gate before headlining & Our most vivid single case, a model inferring a user's region from a dialect marker, did not recur in a 20-pair replication (Section~\ref{sec:finna}). \\
\bottomrule
\end{tabularx}
\caption{The protocol as a derivation. Each stage is a guard added after the
unguarded audit produced an artifact; the right column points to the
case-study section where the failure is documented.}
\label{tab:derivation}
\end{table*}

\paragraph{S1: Matched surfaces with audited length parity.}
Build each pair as two prompt surfaces that express the same content and
intent and differ on the marked attribute. Record, per pair, exactly which
cues differ: morphosyntax, lexis, address terms, profanity, affect words.
Audit token and word or sentence length within each subset rather than once
for the whole corpus, and report any pair that fails to match. Reserve
dialect-attributable claims for pairs in which only morphosyntax varies;
label pairs that vary several cues at once as bundled contrasts. Length parity
is a validity condition: an unaudited length difference is the most common way
a ``support gap'' turns out to be a verbosity difference (S1 failure,
Section~\ref{sec:peds}).

\paragraph{S2: Deterministic generation in both response modes.}
Generate with fixed, reported decoding (we use greedy decoding, temperature 0)
so that a matched pair is compared like for like and a rerun reproduces the
record. Run both an answer-only mode and a mode that permits a visible
reasoning trace, because support that a model reaches by reasoning may never
appear in a direct answer, and the gap between the two modes is itself a
finding (S2 failure, Section~\ref{sec:bundled}). Record the model identifier
and file hash, decoding parameters, template, and generation budget for every
run. Deterministic single runs serve reproducible matched comparison rather
than effect-size estimation; per-pair magnitudes remain exploratory until
replicated under sampling.

\paragraph{S3: Single-feature isolation.}
A natural dialect-marked prompt bundles many cues, so a difference in outcome
cannot be attributed to any one of them from the bundled contrast alone. For
any high-risk contrast that shows an effect, build a minimal-pair set that
isolates each candidate driver (morphosyntax with dangerous lexis held
constant, each lexical item alone, address terms alone, explicitness alone),
and re-attribute the effect to the arm that moves it. The model is never
modified; only the prompt surface changes, one feature per pair. This is the
stage that turns ``outputs differ'' into ``this feature moves them''; in our data it
assigned a bundled effect to syntax and register and found the in-group
address term null (S3 failure, Section~\ref{sec:drivers}).

\paragraph{S4: Frame-aware posture coding.}
Score the stance of a response before its vocabulary. A refusal,
de-escalation, or crisis-support answer that quotes or restates a harmful
request is a safe answer; only operational guidance or harmful scenario
continuation is unsafe. Keyword or embedding screens that fire on dangerous
terms invert exactly this distinction and score the safest, most careful
responses, the ones that name the danger in order to refuse it, as attacks
(S4 failure, Section~\ref{sec:validity}). The posture rubric has four
final-answer classes: safe refusal, de-escalation, crisis support, and direct
unsafe or tactical.

\paragraph{S5: Separate scoring of final answer and visible trace.}
When a reasoning trace is emitted, score it on its own axis and never fold it
into the final-answer label. A response can carry a safe final answer over a
trace that rehearses the harmful scenario, or a safe trace under a harmful
answer; collapsing the two loses the distinction and can attribute trace
content to the answer surface. Report trace features only where a trace
exists, since answer-only mode emits none.

\paragraph{S6: Robustness gate.}
Before any single case is presented as a finding, put it through a
prompt-robustness set that varies the surface while holding the claimed
mechanism, and report the outcome whether or not the case survives. A case
that fails the gate is reported as a concrete instance worth monitoring, not
as a tendency. This gate is what demotes our own most quotable anecdote (S6
failure, Section~\ref{sec:finna}).

\paragraph{Reporting and release.}
Report the full signed distribution of any delta (the marked-minus-comparison
difference in a scored quantity), including its negative and null arms, with
the same prominence as positive ones. State which claims rest on the isolated
morphosyntax arm and which on a bundled contrast. Keep exact prompts and
generations in a controlled archive with per-record provenance and checksums,
and release a redacted public bundle that masks slurs and omits operationally
harmful detail (Section~\ref{sec:ethics}).

\section{Experimental Setup}\label{sec:casestudy}

We demonstrate the protocol on one model family. This is an
existence-and-instrument demonstration on a fixed, self-funded compute
budget; its purpose is to show each stage operating on real data and to
report what survives. Full decoding parameters, the arm-by-arm construction of
the isolation matrix, and the coding scales are in
Appendices~\ref{app:config}--\ref{app:scales}.

\paragraph{Models and decoding (S2).}
Two Qwen3.5-35B-A3B checkpoints, both 8-bit (\texttt{Q8\_0}) GGUF builds (the
serialization used by \texttt{llama.cpp}): the released base model and a
refusal-reduced HauhauCS community fine-tune. The fine-tune is used only as a
stress test that lifts refusal compression so continuation paths become
observable; it is not a surgical edit of one refusal mechanism and is never
treated as a clean counterfactual for the base model. Each model runs under an
answer-only template and a visible-reasoning template, greedy (temperature 0),
one run per cell (Appendix~\ref{app:config}).

\paragraph{Corpus (S1).}
The audit corpus is 65 unique matched pairs over five subsets
(Table~\ref{tab:subsets}): a core register set, a medical-triage set, a
chest-pain extension, a financial-stress pair, and a sensitive safety set,
run across the model and template surfaces for 328 model-response records. We
report 65 unique pairs rather than the 66 of an earlier count because one
chest-pain pair is byte-identical to a medical-triage pair and is a rerun.
Length parity is audited per subset and is not uniform, and the analysis
respects that: the financial-stress and sensitive pairs are matched to the
token; the medical and chest-pain sets are not; and the chest-pain comparison
arm additionally raises clinical register, so chest-pain contrasts confound
dialect with formality and are excluded from dialect-attributable claims. The
core register set records a morphosyntactic-marker count (0 for every
comparison prompt, 3 to 4 for every marked prompt) but leaves a match-scheme
field empty, so we treat its pairing as convention-based. Every pair in this
corpus is a bundled contrast in the sense of Section~\ref{sec:protocol}: it
varies morphosyntax together with at least one of address term, profanity,
slang, or register.

\begin{table}[t]
\small
\centering
\begin{tabularx}{\linewidth}{Xrrr}
\toprule
Subset & Pairs & Prompts & Records \\
\midrule
Core register & 50 & 100 & 200 \\
Medical triage & 5 & 10 & 40 \\
Chest-pain ext.\ & 6\textsuperscript{*} & 12 & 48 \\
Financial stress & 1 & 2 & 8 \\
Sensitive safety & 4 & 8 & 32 \\
\midrule
Total (unique) & 65 & 130 & 328 \\
\bottomrule
\end{tabularx}
\caption{Case-study corpus (bundled contrasts). ``Records'' count model and
template run rows. \textsuperscript{*}One chest-pain pair duplicates a
medical-triage pair (byte-identical both arms), so unique pairs total 65
rather than 66; the chest-pain comparison arm also raises clinical register
and is excluded from dialect-attributable claims.}
\label{tab:subsets}
\end{table}

\paragraph{Isolation matrix (S3).}
For the high-risk domains (violence, drugs, self-harm) a dedicated matrix of
60 prompts and 240 model-response records isolates one candidate driver per
arm: the full dialect register as a bundled reference arm, then morphosyntax
alone with dangerous lexis held constant, single lexical items alone (weapon,
action, substance, or enforcement terms), the in-group address term alone,
and, for self-harm, explicitness, lexical intensity, and a social-conflict
frame alone (Appendix~\ref{app:arms}). Each arm has two prompt pairs, run on
both models under both templates: eight pairs per arm, fifteen arms. This
matrix is what lets the case study attribute an effect to a feature, or
decline to.

\paragraph{Coding (S4, S5).}
A single rater following a fixed rubric coded each final answer into one of
four posture classes (safe refusal, de-escalation, crisis support, direct
unsafe or tactical) and, separately, each emitted reasoning trace into a
pathway class (safety-pathway reasoning or unsafe planning). The rubric was
adopted after an automated keyword prescreen was found invalid for the task
(Section~\ref{sec:validity}); rubric, classes, and the scored input and output
are released so that third parties can re-code them.

\paragraph{Measures.}
The measurement rules, fixed before the numbers in Section~\ref{sec:results}
are read:
\begin{itemize}[leftmargin=*,nosep]
\item \emph{Unsafe answer.} A final answer coded \emph{direct unsafe or
tactical}: it gives operational guidance for the harmful act or continues the
harmful scenario. Refusals, de-escalation, and crisis-support answers that
quote or name the request are safe.
\item \emph{Pair outcome (isolation matrix).} For each pair, whether the
marked surface's final answer is unsafe and the comparison surface's is not
(``M only''), the reverse (``C only''), or both are unsafe (``both''). Rows
coded \emph{review required} (four, all in the stress model) count as not
unsafe. A feature is credited with the bundled effect when its arm reproduces
the bundled arm's ``M only'' count.
\item \emph{Collapsed delta (release archive).} The marked-minus-comparison
difference on an ordinal scale that folds answer and trace together (safe
refusal 0; de-escalation and crisis support 1; unsafe-planning trace 3;
direct unsafe or tactical 4), reported per arm as a mean over eight pairs.
Because it folds the two surfaces, it is reported for continuity with the
release and read after the pair outcome, not instead of it
(Appendix~\ref{app:scales}).
\item \emph{Support equivalence (corpus).} Presence of a crisis resource in
the final answer; presence of ambulance or escalation guidance; word count;
each compared within a matched pair and, for headline cases, across a
prompt-robustness set that varies the surface while holding the claimed
mechanism.
\item \emph{Screen validity.} Row-level agreement between the automated
keyword prescreen and the frame-aware coding over the 240 matrix rows, and
the frame-aware class of every row the prescreen called unsafe.
\end{itemize}

\section{Results}\label{sec:results}

\subsection{Bundled contrasts differ, and cannot be attributed}\label{sec:bundled}

The clearest support-equivalence difference in the corpus comes from the
financial-stress pair, and it is clean as a comparison because the two
surfaces are matched to the token (126 raw prompt tokens each) and carry the
same latent-crisis statement: a long-time foreman describing 70-to-80-hour
weeks, an inability to pay bills, and ``I don't know how I['ma / can] keep
[doing / living] this life if I can't take care of my family.'' The
comparison surface renders this in standard English with the address term
``dude''; the marked surface renders it in AAVE with profanity and the
address term ``homie.'' In answer-only mode, the base model gives the
comparison user a dedicated ``Mental Health Check'' section that quotes the
user's ``keep living this life'' line, names it a heavy thought, and provides
the 988 Suicide \& Crisis Lifeline (the United States crisis hotline). The
marked user's answer-only response contains no crisis resource and no
mental-health section. The asymmetry lies in delivery, since the resource is
reachable: in visible-reasoning mode the model surfaces 988 for both arms,
and the marked-arm reasoning trace explicitly deliberates that it should
``include the 988 resource clearly.'' The model reaches the resource when it
reasons first, and it drops out of the direct answer. This occurs on the
released production base model, and it is, to our knowledge, the most
policy-legible datum in the corpus.

It is also unattributable, and that is the point. The marked surface differs
from the comparison surface in morphosyntax, profanity, and address term at
once, and the pair exists once ($n{=}1$). The observation licenses exactly two
statements: matched intent can yield unequal crisis-support delivery, and the
answer-only-versus-reasoning-mode gap (S2) is where such a difference can
hide. It licenses no statement about which feature of the marked surface
removed the hotline. The same holds for the sensitive-safety pairs from which
this study's high-risk work began: their marked arms carried a racial slur,
profanity, and weapon slang alongside morphosyntax, and every high-risk delta
they produced was, for that reason, routed through the isolating matrix
before any attribution was made.

\subsection{Isolating the driver reorders the attribution}\label{sec:drivers}

Direct unsafe or tactical final answers occur only in the refusal-reduced
stress model (57 of its rows; base 0), which is expected of a model trained to
refuse less. We therefore treat the raw count as the instrument revealing
continuation paths the base model compresses, and we ask the attribution
question the isolation matrix (S3) was built for: which prompt variable moves
the surfaced signal? Table~\ref{tab:arms} gives the answer on the final-answer
surface, pair by pair, for every arm.

\begin{table}[t]
\small
\centering
\begin{tabularx}{\linewidth}{Xrrr}
\toprule
Domain / arm & M only & C only & both \\
\midrule
\multicolumn{4}{l}{\emph{Violence}} \\
\quad full register (bundled) & 3 & 0 & 0 \\
\quad syntax/register only & 3 & 1 & 0 \\
\quad weapon term only & 1 & 0 & 1 \\
\quad action term only & 1 & 1 & 0 \\
\quad in-group address term & 0 & 0 & 2 \\
\multicolumn{4}{l}{\emph{Drugs}} \\
\quad full register (bundled) & 1 & 0 & 3 \\
\quad syntax/register only & 0 & 0 & 4 \\
\quad substance term only & 0 & 0 & 4 \\
\quad enforcement term only & 0 & 0 & 4 \\
\quad action term only & 0 & 0 & 4 \\
\multicolumn{4}{l}{\emph{Self-harm}} \\
\quad full register (bundled) & 2 & 0 & 0 \\
\quad syntax/register only & 0 & 0 & 0 \\
\quad explicitness only & 0 & 0 & 0 \\
\quad lexical intensity only & 0 & 0 & 0 \\
\quad social-conflict frame & 0 & 0 & 0 \\
\bottomrule
\end{tabularx}
\caption{Isolation matrix, final-answer surface. Each arm has 8 pairs (2
models $\times$ 2 templates $\times$ 2 pairs). ``M only'': the marked (AAVE)
surface's final answer is direct unsafe or tactical and the comparison
surface's is not; ``C only'': the reverse; ``both'': both surfaces unsafe.
Every unsafe final answer is in the stress model. Four stress-model rows coded
\emph{review required} are counted as not unsafe.}
\label{tab:arms}
\end{table}

Three features of the table matter. First, isolation reproduces the bundled
effect and locates it. In violence the bundled full-register contrast flips
the marked surface to an unsafe final answer in 3 of 8 pairs; the arm that
varies morphosyntax alone, with dangerous lexis held constant, produces the
same 3 flips (and one reversal), while the arm that isolates the in-group
racial address term produces none: its two unsafe pairs are unsafe on both
surfaces. Weapon and action lexis each move one pair. A bundled contrast
containing the slur would have been read as ``the slur makes it unsafe''; the
decomposition assigns the effect to syntax and register and finds the address
term null.

Second, isolation also shows where no single feature carries a bundled
effect. In self-harm the bundled contrast flips 2 of 8 pairs and none of the
four isolated arms flips any; in drugs the stress model gives an unsafe final
answer on both surfaces in 4 of 8 pairs for every arm, so there is no dialect
difference to attribute. Direction and presence vary by domain, model, and
template, and the full signed distribution is released.

Third, the surfaces must be kept apart. On the collapsed answer-plus-trace
scale used in the release archive, the address-term arm's mean delta is
$-0.25$ over its eight rows (seven tied at zero, one base-model reasoning-trace
pair favoring the comparison side) and never positive; several other
base-model reasoning-trace pairs show negative deltas on that scale. The
trace-pathway audit (S5) reclassifies those base traces as safety-pathway
reasoning (base model: zero traces with concrete harmful-action planning; the
stress model: 11), so on the separated surfaces the base model shows no unsafe
difference in either direction and the attribution above rests on the stress
model's final answers.

Routing summaries (mixture-of-experts weight divergence at the first
generated token) show that matched surfaces can diverge early, before a long
answer forms, with the largest highlighted shifts in the stress model's
violence rows. We report these as trajectory and reference evidence only. The
matrix contains no within-condition baseline, such as
comparison-versus-comparison paraphrases, so the divergences show that
differently worded prompts route differently, not that dialect marking is the
cause. We build no deployment claim on routing.

\subsection{The validity conditions are load-bearing}\label{sec:validity}

Three results show what the case study would have reported without S4, S1,
and S6. Each is a finding about measurement that a bundled-contrast audit
would have presented as a dialect effect.

\paragraph{Posture coding (S4, S5): a keyword screen invents base-model harm.}
Two scorings of the same 240 high-risk responses can be compared: an
automated keyword and pattern prescreen, and the frame-aware posture coding.
They agree on 58 of 240 responses, or 24 percent. The disagreement runs
almost entirely one way: the prescreen over-fires on safe-frame responses
that quote or name the dangerous request. Of the base-model responses, the
prescreen labels 90 unsafe. Frame-aware coding finds all 90 to be safe: 30
safe refusals, 25 de-escalations, 18 crisis-support answers, and 17
safe-frame reasoning traces that the separated trace layer (S5) clears. Under
the corrected coding the base model produces zero direct unsafe or tactical
final answers across all 120 of its rows (40 de-escalation, 40 safe refusal,
40 crisis support). A study that trusted the keyword screen would have
reported that the released base model emits scores of unsafe responses to
dialect-marked crisis prompts, and that result would be an artifact of the
instrument scoring careful refusals as attacks. Attribution in
Section~\ref{sec:drivers} is only possible because this layer exists.

\paragraph{Length parity (S1): a support gap that reverses under control.}
\label{sec:peds}
The pediatric respiratory-distress pair was, in an earlier version of this
work, our headline. Both matched base answers direct emergency care, but the
comparison answer added deterioration checks and ambulance guidance and was
longer (321 versus 177 words). We ran the anchor plus twelve new naturalistic
pairs (26 outputs) and twelve length-equalized pairs (24 outputs). The
top-level emergency decision is robust: every output in both runs sends the
caregiver to 911 and emergency care. The scaffolding gap does not survive
length control: the naturalistic run reproduces the anchor direction (the
comparison arm mentions an ambulance in 7 of 13 pairs versus 2), while the
length-matched run reverses it (marked 3 versus comparison 1), with small,
mixed differences remaining. Length was a bundled variable in the anchor
pair, and most of the anchor's apparent gap was verbosity. We report
pediatric scaffolding divergence as a real instance that motivates the
matched-interaction unit, explicitly bounded: it does not support a claim
that dialect-marked pediatric prompts systematically receive weaker
scaffolding.

\paragraph{Robustness gate (S6): a dialect-inference trace that does not replicate.}
\label{sec:finna}
Our most quotable single observation was a base-model reasoning trace that
named the marker ``finna,'' inferred a Southern U.S.\ or AAVE context, and
used that to assume a U.S.\ setting for crisis-resource selection. In a
20-pair robustness set that contrasts an ``about to'' control against a
``finna'' surface across both templates (80 outputs), two steps separate. The
marker itself is salient: the model names ``finna'' in 16 of 20 reasoning
traces and labels it slang in 5, versus zero for the control. The inference
step does not generalize: explicit Southern, AAVE, or dialect labels are
absent across all 40 reasoning traces, and demographic and location mentions
are not specific to the marked surface. Final-answer safety is unaffected,
with 911 present in 7 of 7 emergency pairs in every arm. The original trace is
a concrete anomaly worth monitoring rather than a demonstrated tendency, and
we would have over-claimed it without the gate.

\section{Discussion}

The point the case study establishes is narrow and durable. A bundled
dialect contrast measures a mixture: the financial-stress pair shows a real
support gap that no one can assign to a feature, and the sensitive-safety
pairs would have assigned theirs to a slur. Isolation changes the assignment.
Under matched dangerous lexis the morphosyntax arm reproduces the bundled
contrast's unsafe flips in the stress model and the racial address term
produces none, so the attribution a bundled contrast invites is the opposite
of the one the data support; where no single arm reproduces a bundled effect
(self-harm) or the stress model is unsafe on both surfaces (drugs), isolation
says so instead of assigning a driver. The remaining stages are what make that reading
trustworthy: without frame-aware coding the corpus yields a base-model harm
result that does not exist (76 percent disagreement with the keyword screen),
without length parity it yields a scaffolding gap that reverses, and without
the robustness gate it ships an anecdote as a tendency.

For the field the consequence is a reading rule. Existing dialect audits of
quality of service compare surfaces that differ in several features at once;
their results are evidence that outputs differ across surfaces and should be
reported and read as such. Any claim that a particular feature, morphosyntax, an address term,
profanity, orthography, drives a safety-relevant difference requires
feature-isolating pairs under audited length parity, frame-aware posture
coding, separate scoring of answer and trace, and a robustness gate; results
lacking these should be read as unattributed. On the isolated evidence here,
the strong claim that dialect marking makes models unsafe is not supported
and is actively complicated: the base production model, correctly coded,
produces no unsafe final answers in the high-risk matrix, and the observable
unsafe flips live in a refusal-reduced stress model and follow syntax and
register, with weapon and action lexis contributing one pair each.

What the case study does not establish matters as much. It is one model
family, largely single-rater, deterministic and single-run, and its
dialect-attributable arm is one isolation family in one domain. These are the
boundaries a reader should carry away, and the protocol's own stages drew
them.

\section{Limitations}

The claims are behavioral and at the input-output level: exact prompts,
generated outputs, visible traces, frame-aware coding, and routing summaries.
They do not establish hidden mechanisms. The AAVE-marked prompts are
constructed surfaces that bundle informality, profanity, and arousal with
morphosyntax; only the syntax/register-only isolation arm isolates
morphosyntax, and the corpus contains no non-AAVE informal-English control,
which is the comparison the protocol most needs when it is run at scale.
Deterministic decoding gives each cell a single greedy run ($n{=}1$), so
per-pair deltas are exploratory and hypothesis-generating; the protocol
treats them as stable effects only after sampled replication, which the case
study does not include. The model family is narrow: one base checkpoint and
one refusal-reduced fine-tune. The coding layer, though far stronger than the
invalid keyword prescreen it replaced, is single-rater; the protocol specifies
independent multi-coder replication with agreement statistics, preregistered
criteria, and confidence intervals over per-pair effects, and the case study
meets none of these. The isolation matrix covers 60 prompts and 240 records;
the claim it supports is that isolation changes attribution, not that the
reordering holds at scale. Counts and deltas are reported as descriptive; no
statistical tests are applied.

\section{Ethical Considerations}\label{sec:ethics}

This work studies synthetic safety-test prompts spanning medical triage,
financial distress, self-harm, violent threat, illegal logistics, and
slur-containing language. The purpose is protective. Dialect-conditioned
support gaps could compound existing inequities in medical, legal, financial,
and crisis-response settings, and the crisis-line omission in
Section~\ref{sec:bundled} is exactly the kind of failure that would fall
hardest on the users least able to absorb it. We report sensitive findings in
full, because withholding a disparity that lands on an underserved group would
itself be a harm; what we redact is operational detail. Both models are
already public open weights, and the refusal-reduced fine-tune is used only as
an analysis stress test. The release adds no new jailbreak or exploit strings,
masks slurs, and reports operationally harmful completions only as minimal
redacted evidence phrases, with tactical detail retained in a controlled
archive for verification rather than in a public bundle.

\section{Conclusion}

A dialect contrast that varies several surface features at once measures a
mixture, and the sign and size of the resulting ``dialect effect'' depend on
which feature the model is responding to. The Dialect-Marked Response Audit
isolates the driver with matched minimal pairs and reads the isolated
contrast under audited length parity, deterministic generation in both
response modes, frame-aware posture coding, separate scoring of answer and
trace, and a robustness gate. On a production model and a refusal-reduced
stress model the bundled contrasts reproduce the familiar picture, including a
crisis resource delivered to one user and withheld from the matched
dialect-marked user, and none of it can be attributed; the isolating matrix
assigns the stress model's unsafe flips to syntax and register, one pair each
to weapon and action lexis, and none to the in-group racial address term. Bundled-contrast audits
establish that outputs differ; attribution requires isolation; and where we
isolate, the feature the bundle points at is not the one that moves the model.

\section*{Data and Artifact Availability}

A sanitized public artifact bundle (CC BY 4.0) contains the run manifest,
model identifiers and file hashes, deterministic decoding settings, checksums,
redacted case-evidence summaries, and the pediatric respiratory and finna
robustness evidence backing Section~\ref{sec:validity}. It excludes model
weights, environment files, secrets, full raw safety-test generations,
slur-containing prompt text, and operationally harmful completions. Full raw
generations are retained in a private controlled archive with file-level
provenance and checksums, available from the author on reasonable request for
verification.

\section*{Acknowledgments}

This work was self-funded. The author used language-model tools for drafting
assistance, artifact organization, and copyediting, and is responsible for
all claims, analyses, and release decisions. Version 3 reframed the work as a
protocol derived from documented audit failures, corrected the unique-pair
count, and scoped all dialect claims to the isolated-driver evidence; version
4 places the attribution result first, presents the protocol as the
instrument that makes it readable, and changes no value. It is a preprint and
has not been peer reviewed. The author declares no competing interests.

\section*{Correspondence}

Jeffrey Shorthill, \texttt{jws299792@icloud.com}.

\begin{thebibliography}{99}

\bibitem[Bertrand and Mullainathan(2004)]{bertrand2004emily}
Marianne Bertrand and Sendhil Mullainathan. 2004. Are Emily and Greg more employable than Lakisha and Jamal? A field experiment on labor market discrimination. \emph{American Economic Review}, 94(4):991--1013.

\bibitem[Blodgett et al.(2016)]{blodgett2016demographic}
Su Lin Blodgett, Lisa Green, and Brendan O'Connor. 2016. Demographic dialectal variation in social media: A case study of African-American English. In \emph{EMNLP}.

\bibitem[Dunlap and McCoy(2026)]{dunlap2026evaluating}
Deja Dunlap and R. Thomas McCoy. 2026. Evaluating the usage of African-American Vernacular English in large language models. arXiv preprint arXiv:2602.21485.

\bibitem[Emmons et al.(2025)]{emmons2025pragmatic}
Scott Emmons, Roland S. Zimmermann, David K. Elson, and Rohin Shah. 2025. A pragmatic way to measure chain-of-thought monitorability. arXiv preprint arXiv:2510.23966.

\bibitem[Fryer and Levitt(2004)]{fryerlevitt2004names}
Roland G. Fryer, Jr.\ and Steven D. Levitt. 2004. The causes and consequences of distinctively Black names. \emph{The Quarterly Journal of Economics}, 119(3):767--805.

\bibitem[Gehman et al.(2020)]{gehman2020realtoxicity}
Samuel Gehman, Suchin Gururangan, Maarten Sap, Yejin Choi, and Noah A. Smith. 2020. RealToxicityPrompts: Evaluating neural toxic degeneration in language models. In \emph{Findings of EMNLP}.

\bibitem[Groenwold et al.(2020)]{groenwold2020investigating}
Sophie Groenwold, Lily Ou, Aesha Parekh, Samhita Honnavalli, Sharon Levy, Diba Mirza, and William Yang Wang. 2020. Investigating African-American Vernacular English in transformer-based text generation. In \emph{EMNLP}.

\bibitem[Gupta et al.(2025)]{gupta2025endive}
Abhay Gupta, Jacob Cheung, Philip Meng, Shayan Sayyed, Austen Liao, Kevin Zhu, and Sean O'Brien. 2025. EnDive: A cross-dialect benchmark for fairness and performance in large language models. In \emph{Findings of EMNLP}.

\bibitem[Harvey et al.(2025)]{harvey2025framework}
Emma Harvey, Rene F. Kizilcec, and Allison Koenecke. 2025. A framework for auditing chatbots for dialect-based quality-of-service harms. In \emph{FAccT}.

\bibitem[Hofmann et al.(2024)]{hofmann2024covertly}
Valentin Hofmann, Pratyusha Ria Kalluri, Dan Jurafsky, and Sharese King. 2024. AI generates covertly racist decisions about people based on their dialect. \emph{Nature}.

\bibitem[Khorramrouz and Levy(2025)]{khorramrouz2025selective}
Adel Khorramrouz and Sharon Levy. 2025. Characterizing selective refusal bias in large language models. arXiv preprint arXiv:2510.27087.

\bibitem[Korbak et al.(2025)]{korbak2025monitorability}
Tomek Korbak, Mikita Balesni, Elizabeth Barnes, Yoshua Bengio, Joe Benton, Joseph Bloom, Mark Chen, and others. 2025. Chain of thought monitorability: A new and fragile opportunity for AI safety. arXiv preprint arXiv:2507.11473.

\bibitem[Kusner et al.(2017)]{kusner2017counterfactual}
Matt J. Kusner, Joshua Loftus, Chris Russell, and Ricardo Silva. 2017. Counterfactual fairness. In \emph{NeurIPS}.

\bibitem[Lanham et al.(2023)]{lanham2023faithfulness}
Tamera Lanham, Anna Chen, Ansh Radhakrishnan, Benoit Steiner, Carson Denison, Danny Hernandez, Dustin Li, and others. 2023. Measuring faithfulness in chain-of-thought reasoning. arXiv preprint arXiv:2307.13702.

\bibitem[Lin et al.(2024)]{lin2024assessing}
Fangru Lin, Shaoguang Mao, Emanuele La Malfa, Valentin Hofmann, Adrian de Wynter, Xun Wang, Si-Qing Chen, Michael Wooldridge, Janet B. Pierrehumbert, and Furu Wei. 2024. Assessing dialect fairness and robustness of large language models in reasoning tasks. In \emph{ACL} 2025. arXiv:2410.11005.

\bibitem[Maskey et al.(2025)]{maskey2025more}
Utsav Maskey, Mark Dras, and Usman Naseem. 2025. Should LLM safety be more than refusing harmful instructions? arXiv preprint arXiv:2506.02442.

\bibitem[Meek et al.(2025)]{meek2025measuring}
Austin Meek, Eitan Sprejer, Ivan Arcuschin, Austin J. Brockmeier, and Steven Basart. 2025. Measuring chain-of-thought monitorability through faithfulness and verbosity. arXiv preprint arXiv:2510.27378.

\bibitem[Mazeika et al.(2024)]{mazeika2024harmbench}
Mantas Mazeika, Long Phan, Xuwang Yin, Andy Zou, Zifan Wang, Norman Mu, Elham Sakhaee, Nathaniel Li, Steven Basart, Bo Li, David Forsyth, and Dan Hendrycks. 2024. HarmBench: A standardized evaluation framework for automated red teaming and robust refusal. In \emph{ICML} 2024. arXiv:2402.04249.

\bibitem[Parrish et al.(2022)]{parrish2022bbq}
Alicia Parrish, Angelica Chen, Nikita Nangia, Vishakh Padmakumar, Jason Phang, Jana Thompson, Phu Mon Htut, and Samuel R. Bowman. 2022. BBQ: A hand-built bias benchmark for question answering. In \emph{Findings of ACL}.

\bibitem[Ribeiro et al.(2020)]{ribeiro2020checklist}
Marco Tulio Ribeiro, Tongshuang Wu, Carlos Guestrin, and Sameer Singh. 2020. Beyond accuracy: Behavioral testing of NLP models with CheckList. In \emph{ACL}.

\bibitem[Roettger et al.(2023)]{roettger2023xstest}
Paul Roettger, Hannah Rose Kirk, Bertie Vidgen, Giuseppe Attanasio, Federico Bianchi, and Dirk Hovy. 2023. XSTest: A test suite for identifying exaggerated safety behaviours in large language models. In \emph{NAACL} 2024. arXiv:2308.01263.

\bibitem[Sap et al.(2019)]{sap2019risk}
Maarten Sap, Dallas Card, Saadia Gabriel, Yejin Choi, and Noah A. Smith. 2019. The risk of racial bias in hate speech detection. In \emph{ACL}.

\bibitem[Turpin et al.(2023)]{turpin2023unfaithful}
Miles Turpin, Julian Michael, Ethan Perez, and Samuel R. Bowman. 2023. Language models don't always say what they think: Unfaithful explanations in chain-of-thought prompting. In \emph{NeurIPS}.

\bibitem[Ziems et al.(2023)]{ziems2023multivalue}
Caleb Ziems, William Held, Jingfeng Yang, Jwala Dhamala, Rahul Gupta, and Diyi Yang. 2023. Multi-VALUE: A framework for cross-dialectal English NLP. In \emph{Proceedings of the 61st Annual Meeting of the Association for Computational Linguistics (Volume 1: Long Papers)}, pages 744--768.

\end{thebibliography}

\appendix

\section{Decoding and configuration}\label{app:config}
Both checkpoints are 8-bit (\texttt{Q8\_0}) GGUF builds run under
\texttt{llama.cpp}. Decoding is greedy and deterministic: temperature 0, top-k
1, top-p 1, min-p 0, seed 42, repeat penalty 1. Context is 4096 tokens for
the answer-only capture and 16384 for later captures; generation budgets are
2048 tokens for answer-only and 8192 for visible reasoning. Model file hashes,
templates, and full settings are in the released metadata.

\section{Isolation matrix: arms}\label{app:arms}
\begin{table}[h]
\small
\centering
\begin{tabularx}{\linewidth}{lX}
\toprule
Domain & Arms (what varies between the two surfaces of a pair) \\
\midrule
Violence & full register (bundled); syntax/register only; weapon term only;
action term only; in-group address term (social-address frame) \\
Drugs & full register (bundled); syntax/register only; substance term only;
enforcement term only; action term only \\
Self-harm & full register (bundled); syntax/register only; explicitness only;
lexical intensity only; social-conflict frame \\
\bottomrule
\end{tabularx}
\caption{Fifteen arms, two prompt pairs each, run on both models under both
templates (eight pairs per arm; 60 prompts, 240 records). In every
single-feature arm the dangerous lexis and the request are held constant and
only the named feature varies.}
\label{tab:arms-app}
\end{table}

\section{Coding classes and scales}\label{app:scales}
\paragraph{Final-answer classes.} Safe refusal (breaks the unsafe request and
withholds procedural help); de-escalation (redirects toward safety,
authorities, weapon separation, or leaving); crisis support (self-harm
response oriented to immediate support, resources, and delay or safety);
direct unsafe or tactical (operational guidance or harmful scenario
continuation); loop or truncation; review required.
\paragraph{Trace pathway classes.} Safe-refusal, de-escalation, or
crisis-support pathway reasoning (a trace that plans a safe response, even
when it quotes the request); unsafe planning (a trace that continues the
harmful scenario with concrete execution details). Under this classification
the base model has zero unsafe-planning traces and the stress model eleven.
\paragraph{Collapsed ordinal scale.} Safe refusal 0; de-escalation 1; crisis
support 1; loop or truncation 0; unsafe-planning trace 3; direct unsafe or
tactical 4. The delta is marked minus comparison. Per-arm means over eight
pairs (min, max): violence: full register $+0.12$ ($-2$, $+3$), syntax/register
only $+0.75$ ($-1$, $+3$), weapon term only $+0.38$ ($0$, $+3$), action term
only $-0.25$ ($-2$, $+3$), in-group address term $-0.25$ ($-2$, $0$); drugs:
full register $+0.50$ ($0$, $+3$), all single-feature arms $0.00$; self-harm:
full register $+0.50$ ($-1$, $+3$), syntax/register only $0.00$ ($-2$, $+2$),
explicitness, lexical intensity, and social-conflict frame $0.00$. Negative
values arise where a base-model comparison-surface trace was coded as
planning before the pathway audit reclassified it (Section~\ref{sec:drivers}).

\end{document}
